\section{Methodology}\label{sec:method}
The study is a local, exploratory development analysis. V7-B and its baselines were shaped while seeing these data, no sealed or independent holdout exists, and the confirmatory gate K5 (at least 60 projects, at least 250 adjudicated divergent pairs, kappa of at least 0.7, clean leakage audit) was not attempted. The real-pair feasibility gate K1 (at least 100 runnable real pairs from at least 8 projects under validated Linux isolation) is neither triggered nor cleared; K1, K5 and all Phase 3 work are deferred. Under these conditions every inferential statement below is descriptive.

The protocols for E3--E7 and the hypothesis and decision rule for E1 were written and committed before the corresponding runs (\path{reports/phase2/E3_E7_PROTOCOLS.md}, \path{reports/phase2/LOCAL_EXPERIMENT_RERUN_MATRIX.md}); this pre-registration is internal and does not make the study confirmatory. Pipeline: run V7-B on the 321 pairs through its command-line interface (per-run limits of 2~s CPU and 5~s wall, four workers in E1, E3 and E4, sequential execution in E6); run the baselines on identical inputs; build one table of (pair, system) rows; analyse with a project-cluster bootstrap and project-level sign-flip tests (Section~\ref{sec:metrics}); verify every evidence directory with \texttt{verify} (E7). A pair is \emph{driven} when the engine supports its entry (top-level function); 9 pairs are not driven (eight class units and one unsupported entry policy) and count as failed or missing runs, that is abstention, in intention-to-treat accounting.
